\noindent\textbf{Objective.} Decoders for intracortical brain--computer interfaces lose accuracy as recordings change.
How fast this happens, why, and which corrections work have not been compared across individuals with one method.

\noindent\textbf{Approach.} We analyzed public recordings from three monkeys and six BrainGate participants, spanning
up to \val{hum:years:max} years, with one pipeline. Ridge decoders trained on one day were tested 1 to 480 days later
after a ladder of corrections, from per-session normalization of each channel and label-free alignment to channel
re-weighting and refitting with labeled trials. An alignment fitted with labels gave a ceiling for any alignment of
the dominant neural subspace. Electrode statistics came from 20 arrays in 14 participants.

\noindent\textbf{Main results.} Decay differed widely between individuals (mixed model, $P = \val{het:all:curves:p}$):
a fixed decoder lost half its accuracy within a day in some and only after \val{T6:thalftuned} days in another. Weak
regularization exaggerated decay. No label-free alignment we tested improved on per-session normalization, although
the ceiling recovered \val{hum:ceil:rot:min}--\val{mon:ceil:rot:max}\% of the loss at 120 to 480 days; even the
ceiling left \val{all:ceil:unrec:min}--\val{all:ceil:unrec:max}\%. In monkeys the ceiling beat channel re-weighting,
pointing to activity remixed across channels; in people re-weighting did better. In people, about 100 labeled trials
restored most accuracy, and a fresh, strongly regularized decoder used them at least as well as a refit toward the old
one. In a pilot, a pretrained foundation decoder also needed labels. Most brief electrode silences could not be told
apart from a permutation null; month-long silences were real but often reversed.

\noindent\textbf{Significance.} Better label-free alignment is possible but cannot remove the need for labels. Decoder
longevity should be compared between tuned decoders, label-free methods against per-session normalization, and
recalibration schedules should be set per person. We release a \val{bench:npairs}-pair benchmark spanning 480 days, a
calibrated simulator and all code.
